% page 408 du fac-simile (image p-407 du scan)
% bloc 149.9 x 246.3 mm ; marge gauche 8.0 mm
\pgc{-12.700mm}{0.000mm}{
\l{\marge[78.592mm]{30.698mm}{\ornamento{13.475mm}{ornements/letroj/litero-O.png}}}
\l{}
\l{\cel{3}400.}
\l{}
\l{}
\l{}
\l{}
\l{}
\l{}
\l{}
\l{}
\l{\cel{3}o\cel{1}\sou{(d)}.\cel{2}Konjunciono qua indikas alternativo.}
\l{}
\l{\sou{oaziso}. Loko sola, kun vejetantaro, en dezerto de sablo.}
\l{\cel{3}DEFIRS.}
\l{}
\l{\sou{obcena}. Qua shokas pudoro.- DEFIS.}
\l{}
\l{\sou{obediar}.\cel{2}(\sou{trans}.) I. Agar konforme a to quon imperas ulo.}
\l{\cel{3}- II. (\sou{aludante kozo}).Subisar\cel{2}la ago da ulu o da ulo.}
\l{\cel{3}- EFIS.}
\l{}
\l{\sou{obelisko}. Monumento (maxim-multa-kaze monolita) taliita}
\l{\cel{3}forme di agulo quadrangula. - DEFIRS.}
\l{}
\l{\sou{obeza}. Tro repleta e, konseque, di qua la ventro (che}
\l{\cel{3}la viri) o la hanchi e la gropo (che la mulieri)}
\l{\cel{3}salias exajere pro akumulaji di graso sub-pele.}
\l{\cel{3}- EFIS.}
\l{}
\l{\sou{objecionar}. (\sou{trans., ad ulu}). I. Opozar (ulo) ad afirmo}
\l{\cel{3}por kombatar olu. - II. Opozar (ulo) a propozo, a pro-}
\l{\cel{3}jeto, por repulsar olu. - EFIS.}
\l{}
\l{\sou{objekto}. I. To quo esas vidata, rigardata.-II. To qua}
\l{\cel{3}su prizentas en la spirito, en la penso. - III.}
\l{\cel{3}(\sou{filoz}.) To quo pensesas (opoze a to qua pensas).}
\l{\cel{3}- IV. (\sou{gram}.) Termino di la propoziciono a qua agas}
\l{\cel{3}la subjekto. V. To a quo iras la sentimenti di la}
\l{\cel{3}anmo, di la psiko.- VI. To a quo iras volo. - DEFIRS.}
\l{}
\l{\sou{oblato}. Pano sen-hefa, tranchita ronde, uzata por siglar}
\l{\cel{3}kuverti. DRS.}
\l{}
\l{\sou{obleo}. Kukajo dina, (normale) volvita forme di cilindro}
\l{\cel{3}vakua, o di korneto. FRS.}
\l{}
\l{\sou{obligar}.\cel{2}(\sou{trans., ad ulo)}.\cel{2}Ligar (ulu) per lego, per kon-}
\l{\cel{3}venciono, quan lu mustas obediar. - DEFIS.}
\l{}
\l{\sou{obligaciono}. (\sou{financo)}. Akto qua reprezentas parto di la}
\l{\cel{3}kapitalo pruntita da societo financala, industriala,}
\l{\cel{3}e c. - qua yurizas pri interesto-procento fixa, e qua}
\l{\cel{3}rimborsesos pos fristo determinita. - DEFIRS.}
\l{}
\l{\sou{obliqua}. I. Qua deviacas de vertikalo. - II. Qua agas per}
\l{\cel{3}moyeni nedireta. - EFIS.}
}
